% Original PDF: source/普通高考/2015/2015天津文.pdf, page 1, question 3.
% pdfimages -list found no embedded bitmap; the flowchart is vector/text line art.
% Grid evidence: tmp/redraw_flowchart_2015_tianjin_liberal/grid/q3_block_grid.jpg;
% comparison source is the no-grid crop q3_original_full.png from that page.
% Node -> directed edges: 开始→S=10,i=0→i=i+1→S=S−i→S≤1?;
% 是→输出S→结束；否→right loop→shared entry above i=i+1.
% All borders and connectors are solid; no dashed segments. Node text is centered
% with parindent=0pt; the shared loop entry has one left arrow and one downward arrow.
\begin{tikzpicture}[
  >=Stealth,
  line width=0.5pt,
  line cap=round,
  line join=round,
  font=\large,
  flowbox/.style={draw,align=center,anchor=center,
    inner xsep=4pt,inner ysep=2pt,
    execute at begin node={\setlength{\parindent}{0pt}}},
  io/.style={flowbox,shape=trapezium,trapezium left angle=75,trapezium right angle=105},
  flowarrow/.style={->,line width=0.5pt},
]
  \node[flowbox,rounded corners=0.18cm,minimum width=1.35cm,minimum height=0.68cm]
    (start) at (0,8.15) {开始};
  \node[flowbox,minimum width=3.15cm,minimum height=0.70cm]
    (init) at (0,6.90) {$S=10,i=0$};
  \node[flowbox,minimum width=2.20cm,minimum height=0.68cm]
    (increment) at (0,5.62) {$i=i+1$};
  \node[flowbox,minimum width=2.25cm,minimum height=0.68cm]
    (update) at (0,4.30) {$S=S-i$};
  \node[flowbox,shape=diamond,aspect=2.55,minimum width=3.80cm,
    minimum height=1.00cm,inner sep=0pt]
    (test) at (0,2.90) {$S\leq1?$};
  \node[io,minimum width=2.30cm,minimum height=0.78cm]
    (output) at (0,1.42) {输出 $S$};
  \node[flowbox,rounded corners=0.18cm,minimum width=1.35cm,minimum height=0.68cm]
    (finish) at (0,0.10) {结束};

  % Main sequence with a single downward arrow at the shared loop entry.
  \draw[flowarrow] (start.south) -- (init.north);
  \coordinate (loopjoin) at (0,6.32);
  \draw (init.south) -- (loopjoin);
  \draw[flowarrow] (loopjoin) -- (increment.north);
  \draw[flowarrow] (increment.south) -- (update.north);
  \draw[flowarrow] (update.south) -- (test.north);
  \draw[flowarrow] (test.south) -- node[right] {是} (output.north);
  \draw[flowarrow] (output.south) -- (finish.north);

  % 否 branch: right, up, then left back into the shared entry above i=i+1.
  \coordinate (right) at (2.75,2.90);
  \coordinate (righttop) at (2.75,6.32);
  \draw (test.east) -- node[above] {否} (right);
  \draw (right) -- (righttop);
  \draw[flowarrow] (righttop) -- (loopjoin);
\end{tikzpicture}
